% Document Setup ---------------------------------------------------------

% Set document class.
\documentclass[10pt]{article}

% Set document details.
\title{\Large\textsc{Frontiers in Computing: Discussion 2}}
\author{\large\textsc{Rento Saijo}}
\date{2026-09-08}

% Load packages.
\usepackage[letterpaper, margin = 1in]{geometry}
\usepackage{amsmath, amssymb}
\usepackage{enumitem}
\usepackage{graphicx}
\usepackage{fancyhdr, lastpage}
\usepackage{tcolorbox}
\usepackage[hidelinks]{hyperref}

% Page Layout ------------------------------------------------------------

% Configure equations and paragraphs.
\allowdisplaybreaks
\setlength{\parindent}{0pt}
\setlength{\parskip}{0.3em}
\clubpenalty = 10000
\widowpenalty = 10000

% Configure page footer.
\pagestyle{fancy}
\fancyhf{}
\fancyfoot[C]{\textsc{Page \thepage\ of \pageref{LastPage}}}
\renewcommand{\headrulewidth}{0pt}

% Helpers ----------------------------------------------------------------

% Define problem environment.
\newtcolorbox{problem}[1][]{
    colframe = darkgray,
    arc      = 0.1mm,
    boxrule  = 2pt,
    boxsep   = 1mm,
    title    = {\textsc{#1}}
}

% Content ----------------------------------------------------------------

% Render document title.
\begin{document}
\maketitle
\thispagestyle{fancy}
\vspace{-1em}\rule{\linewidth}{0.6pt}\vspace{0.5em}

% Introduce Internet infrastructure and protocols.
\section*{The Internet and its protocols}

\begin{itemize}[leftmargin=1.25em, itemsep=0.3em, parsep=0pt, topsep=0pt]
    \item We can understand the Internet through both its components and the services they provide. Hosts, or end systems, include personal computers, phones, servers, and connected appliances. Communication links join these devices to packet switches, including routers and switches, which move information toward its destination. Networks operated by different organizations interconnect to form a network of networks. From an application's perspective, this infrastructure supplies a communication service: programs running on different hosts can exchange messages without implementing every mechanism that carries those messages across the Internet.

    \item Communication requires agreement about what messages mean and when to send them. A protocol specifies message formats, their order, and the actions taken when messages are transmitted or received. Asking someone for the time illustrates the idea: a greeting establishes a conversation, a question requests information, and a response supplies it. Network protocols similarly coordinate exchanges between communicating systems. Common specifications allow equipment and software from different organizations to work together. The Internet Engineering Task Force develops Internet standards and publishes protocol specifications through the Request for Comments, or RFC, series.
\end{itemize}

% Describe end systems, access networks, and physical media.
\section*{The network edge}

\begin{itemize}[leftmargin=1.25em, itemsep=0.3em, parsep=0pt, topsep=0pt]
    \item At the edge, applications run on hosts: clients request services, while servers respond to those requests. An access network connects a host to its first router on a path into the wider Internet. Residential cable, DSL, institutional Ethernet, Wi-Fi, and cellular access provide different ways to make this connection. Two useful questions are how much transmission capacity is available and who shares it. Cable subscribers may compete with neighbors for local capacity; DSL uses a dedicated copper pair to the telephone company's central office, although later parts of the Internet path can still be shared.

    \item A home network commonly combines a modem, router, wired connections, and a wireless access point, sometimes within one device. Institutional networks connect many hosts through switches and routers, while data centers interconnect large collections of servers. These arrangements depend on physical media. Twisted-pair copper and coaxial cable guide electrical signals, optical fiber guides light, and wireless links carry radio signals through the environment. Fiber supports high capacity with low error rates. Wireless access offers mobility, but obstructions, reflection, interference, and fading affect reception; signals accessible to nearby receivers also create opportunities for eavesdropping.

    \item Before transmission, a host divides a large application message into packets and adds protocol information in headers. If a packet contains \(L\) bits and its outgoing link transmits at \(R\) bits per second, placing all its bits onto that link takes \(L/R\) seconds. This relationship connects packet size to link capacity: larger packets take longer to transmit, while a higher rate reduces transmission time. It does not yet describe how long the signal travels through the medium or how long the packet waits elsewhere along its route.
\end{itemize}

% Explain forwarding, switching, and Internet interconnection.
\section*{The network core}

\begin{itemize}[leftmargin=1.25em, itemsep=0.3em, parsep=0pt, topsep=0pt]
    \item Packets traverse the core through routers and their interconnected links. Forwarding is a local operation: a router reads destination information, consults its forwarding table, and selects an outgoing link. Routing determines network paths and supplies information for these tables. In the driving analogy, planning a journey from San Jose to Northampton corresponds to routing; choosing the correct outgoing road at an individual interchange corresponds to forwarding. Repeated local decisions therefore carry packets along paths established through routing.

    \item Packet switching shares links among packets as they arrive. With store-and-forward operation, a router receives an entire packet before transmitting it onto the next link. Packets wait in a buffer when that outgoing link is busy. If arriving traffic exceeds outgoing capacity, the queue grows; a full buffer causes packet loss. Sharing capacity accommodates changing demand efficiently, although competition for resources produces variable delays.

    \item Circuit switching allocates resources to a connection in advance. Frequency-division multiplexing reserves a frequency band, while time-division multiplexing reserves recurring time slots. Reserved capacity remains assigned even during silence. Packet switching uses statistical multiplexing, sharing capacity according to current activity. In the lecture's example, a 1 Gbps link can reserve 100 Mbps each for ten users. If users transmit only 10\% of the time, packet switching can accommodate a larger group because relatively few users transmit simultaneously. Overlapping bursts can still cause congestion.

    \item Connecting every access network directly to every other one would require an impractical number of links. Transit providers connect many networks, regional networks aggregate access networks, and competing providers exchange traffic through peering connections and Internet exchange points. Content providers can operate private backbones connecting their data centers and connect directly to networks nearer their users. Through these overlapping relationships, independently administered networks cooperate to provide connectivity without one organization operating the entire Internet.
\end{itemize}

% Relate delay, congestion, measurements, and throughput.
\section*{Delay, loss, and throughput}

\begin{itemize}[leftmargin=1.25em, itemsep=0.3em, parsep=0pt, topsep=0pt]
    \item Four kinds of delay affect a packet's progress. Processing includes checking the packet and choosing its outgoing link; queuing is time awaiting transmission. Transmission takes \(L/R\) seconds to place its bits onto a link, while propagation is their travel time through the medium. In the tollbooth analogy, releasing ten cars at twelve seconds per car takes two minutes; traveling 100 km at 100 km/h takes each car an hour. The last car reaches the next booth 62 minutes after service begins for the first car. Faster booth service reduces release time without changing road travel time.

    \item Queuing depends strongly on traffic intensity. If packets average \(L\) bits, arrive at an average rate of \(a\) packets per second, and leave over a link of rate \(R\), the ratio \(La/R\) compares incoming demand with capacity. With variable arrivals, queues can grow sharply as this ratio approaches one. Sustained demand above capacity creates an ever-growing backlog in an idealized infinite buffer; real buffers eventually fill and drop packets. Packet loss can lead to retransmission, while mechanisms such as TCP congestion control adjust sending rates in response to network conditions.

    \item Traceroute makes parts of this behavior observable by sending probes toward successive routers and recording round-trip times. Each measurement includes travel toward the responding router and the reply's journey back. Multiple probes can produce different times because their waiting conditions differ, so a later router can even have a smaller recorded round-trip time than an earlier one. Large increases can reflect long-distance links, including an ocean crossing. An unanswered probe leaves a missing measurement; it does not by itself establish that ordinary traffic cannot reach the destination.

    \item Throughput measures the rate at which data reaches a receiver over an interval. An otherwise unconstrained path with two link rates \(R_1\) and \(R_2\) is limited by its bottleneck, \(\min(R_1,R_2)\). Shared capacity also matters: if ten transfers divide a core link's rate \(R\) equally, each gets \(R/10\), which must be compared with its access-link rates. Link capacity, delay, and achieved throughput describe different aspects of performance.
\end{itemize}

% Connect protocol services through layering and encapsulation.
\section*{Layers and encapsulation}

\begin{itemize}[leftmargin=1.25em, itemsep=0.3em, parsep=0pt, topsep=0pt]
    \item Layering organizes networking into manageable responsibilities, with each layer using services below it to provide services above. The airline analogy shows how this organization helps: ticketing, baggage handling, gate operations, and aircraft routing contribute distinct functions to one journey. A change within one function need not require redesigning every other function if their service interfaces remain stable. Networking gains the same modularity, allowing us to reason about application communication without simultaneously describing every detail of cables, radio signals, and router hardware.

    \item The five-layer Internet model begins with the application layer, whose protocols govern exchanges among application components. The transport layer provides process-to-process communication, distinguishing applications running on a host. The network layer provides host-to-host packet delivery across networks. Below it, the link layer transfers data between adjacent devices over an individual link, and the physical layer transmits bits through the medium. The services offered at these boundaries matter: IP supplies best-effort host-to-host delivery, while TCP implements reliable process-to-process transfer using mechanisms above that underlying service, which can lose packets.

    \item Encapsulation carries information between these layers. A transport header accompanies an application message to form a segment; a network header produces a datagram; and link-layer information packages it into a frame for transmission. At the receiver, the corresponding layers interpret and remove this information as the message moves upward. Headers contain information needed at their own layer, such as process identifiers or host addresses. Routers and switches carry out forwarding through lower-layer functions, so they can move application traffic without implementing the application's meaning. Logical communication between peer layers thus rests on actual movement down, across, and up the stack.
\end{itemize}

% Explain attacks and corresponding security protections.
\section*{Networks under attack}

\begin{itemize}[leftmargin=1.25em, itemsep=0.3em, parsep=0pt, topsep=0pt]
    \item Security becomes essential when communicating parties cannot assume mutual trust, as the Internet's early community more readily could. Packet sniffing copies traffic visible on a shared medium, potentially exposing unprotected information; packet-analysis tools such as Wireshark also use this capability for legitimate inspection. Spoofing introduces false source information, making a claimed identity unreliable. Denial-of-service attacks exhaust network or server resources, preventing legitimate users from receiving service. Distributed attacks coordinate traffic from many machines, often compromised hosts, increasing the burden on the target and complicating attempts to filter the traffic.

    \item Different protections address different weaknesses. Authentication establishes identity, encryption protects message contents from disclosure, and digital signatures help verify origin and detect alteration. Access controls determine which authenticated users or devices may perform particular actions. Firewalls apply traffic rules at network boundaries or within networks and can help limit unwanted access and attack traffic. These mechanisms complement one another: proving identity does not itself keep a message confidential or prevent resource exhaustion. Designing security into network architecture therefore means considering identities, information, and shared resources throughout the layers that make communication possible.
\end{itemize}

% End document.
\end{document}
